\documentclass[11pt,letterpaper]{article}
\usepackage[margin=1in]{geometry}
\usepackage[T1]{fontenc}
\usepackage{lmodern,amsmath,amssymb,amsthm,booktabs,microtype}
\usepackage[hidelinks]{hyperref}
\hypersetup{pdftitle={A formal Jacobian ideal certificate for a mixed quartic},pdfsubject={A fixed-quartic Lean certificate with a separate written homogeneous extension}}
\usepackage{listings}
\lstset{basicstyle=\ttfamily\small,breaklines=true,columns=fullflexible}
\newtheorem{theorem}{Theorem}[section]
\newtheorem{lemma}[theorem]{Lemma}
\newtheorem{proposition}[theorem]{Proposition}
\theoremstyle{remark}\newtheorem{remark}[theorem]{Remark}
\DeclareMathOperator{\Der}{Der}
\DeclareMathOperator{\End}{End}
\DeclareMathOperator{\Hess}{Hess}
\DeclareMathOperator{\adj}{adj}
\DeclareMathOperator{\tr}{tr}
\DeclareMathOperator{\Soc}{Soc}
\newcommand{\QQ}{\mathbb Q}
\newcommand{\DeltaA}{\Delta_k(A)}
\title{A formal Jacobian ideal certificate for a mixed quartic}
\author{}
\date{12 October 2026}
\begin{document}
\maketitle
\begin{abstract}
For the explicit ring
\(A=\QQ[x_0,x_1,x_2,x_3]/(\sum_i x_i^4+\sum_{i<j}x_i^2x_j^2)\),
we give a complete Lean-checked construction of a differential operator of order
at most two outside the full algebra generated by multiplication operators and
all derivations. A finite rational coefficient functional detects its failure
to preserve the Jacobian ideal. The proof covers arbitrary finite derivation
words, including cancellation among words of different lengths. We also give
a written extension of the Hessian construction to isolated ordinary
homogeneous hypersurfaces, using the classical complete-intersection socle
theorem. That general extension is not formalized here. The three-variable
operator formula is already in the literature, and publication priority for
the higher-dimensional argument has not been established. This note does not
resolve the unrestricted Nakai conjecture.
\end{abstract}

\section{Scope and prior work}
For a commutative algebra \(A\) over a characteristic-zero field \(k\), let
\[
 \Delta_k(A)=k\langle m_a,\delta : a\in A,\ \delta\in\Der_k(A)\rangle
 \subseteq\End_k(A),\qquad m_a(b)=ab.
\]
This is an algebra with no bound on word length. If
\([T,m_a]=Tm_a-m_aT\), the condition
\begin{equation}\label{eq:order}
 [[[T,m_a],m_b],m_c]=0\qquad(a,b,c\in A)
\end{equation}
is the usual faithful endomorphism definition of differential order at most two.
Our formal statement produces a \(T\) satisfying \eqref{eq:order} and
\(T\notin\Delta_k(A)\).

This distinction matters. Exclusion from the module spanned by derivations and
products of two derivations does not by itself exclude membership in
\(\Delta_k(A)\): longer words could cancel in their highest-order parts.
The introduction and principal theorem of the 8 October 2026 version of
Yau--Zhu--Zuo \cite{YZZ} explicitly use generation order by order and
\(\operatorname{der}^2(A)\ne\Der^2(A)\). We do not replace the unrestricted
algebra by that shorter-word module.

\paragraph{Prior formula and method.}
The operator below is not claimed to be a new formula in three variables.
Diethorn--Jeffries--Miller--Packauskas--Pollitz--Rahmati--Vassiliadou
\cite[Sections 4.1.3 and 4.5, equations (33), (35), Theorem 4.5.1]{DJM}
give precisely \(Q_r/(d-1)^2\) in that case, after conversion from divided
partials and inclusion of the first-order correction. Their introduction
also discusses the earlier work of Vigu\'e. Stable-ideal separation is a
known method; it occurs, for example, in the Fermat-type work of
Brumatti--Veloso \cite{BV}. Hessian-adjugate constructions also overlap
\cite{YZZ}. Traves \cite{Traves} proves the full conjecture for varieties
with smooth normalization. Consequently neither a general novelty claim nor
a claim to have solved a previously unverified singularity class follows from
this note. The formal verification and explicit full-word certificate are the
precisely documented outputs.

\section{The fixed quartic theorem}
Put \(P=\QQ[x_0,x_1,x_2,x_3]\) and
\begin{equation}\label{eq:F}
 F=\sum_{i=0}^3x_i^4+\sum_{0\le i<j\le3}x_i^2x_j^2.
\end{equation}
Let \(G_i=\partial_iF\), \(H=(\partial_i\partial_jF)\),
\(C=\adj(H)\), \(h=\det H\), \(a_i=C_{i0}\), and
\[
 B_{ij}=a_ix_j+x_ia_j-x_0C_{ij}.
\]
Define the explicit polynomial operator
\begin{equation}\label{eq:QP}
 \mathcal Q(p)=\sum_{i,j=0}^3B_{ij}\partial_i\partial_jp
              +6\sum_{i=0}^3a_i\partial_ip.
\end{equation}
The partials are ordinary, not divided. Both ordered off-diagonal terms
occur in the double sum.

\begin{theorem}[Fully formalized fixed case]\label{thm:fixed}
Let \(A=P/(F)\) and let \(\pi:P\to A\) be the quotient map. There is
\(Q\in\End_{\QQ}(A)\) such that
\[
 Q(\pi(p))=\pi(\mathcal Q(p))\quad(p\in P),
 \qquad [[[Q,m_u],m_v],m_w]=0\quad(u,v,w\in A),
 \qquad Q\notin\Delta_{\QQ}(A).
\]
\end{theorem}
The Lean declarations are
\begin{lstlisting}
NakaiQuartic.quartic_full_operator_separation
NakaiExpected.expected_full_operator_separation
\end{lstlisting}
The latter restates the literal ten-term polynomial, multiplication maps,
commutators, explicit polynomial operator, and unrestricted generating set
independently of the implementation abbreviations. Neither statement assumes
the desired nonmembership conclusion. Sections~\ref{sec:stable}--\ref{sec:finite}
provide the mathematical proof and its certificates.

\section{The ideal preserved by every derivation}\label{sec:stable}
\begin{lemma}\label{lem:stable}
Let \(P=k[x_1,\ldots,x_n]\), \(f\in P\), and \(A=P/(f)\).
The image \(J_A\) of the gradient ideal \(J=(f_1,\ldots,f_n)\)
is preserved by every derivation of \(A\). Every element of
\(\Delta_k(A)\) therefore preserves \(J_A\).
\end{lemma}
\begin{proof}
For \(\delta\in\Der_k(A)\), choose representatives \(u_i\in P\) of
\(\delta(\overline{x_i})\), and put
\(\widetilde\delta=\sum_i u_i\partial_i\).
The two induced derivations agree on all polynomial generators, hence agree
on \(P\). In particular \(\widetilde\delta(f)=vf\) for some \(v\in P\).
Differentiating this equality gives
\begin{equation}\label{eq:stability}
 \widetilde\delta(f_j)=v f_j+(\partial_jv)f
                      -\sum_i(\partial_ju_i)f_i.
\end{equation}
Thus \(\widetilde\delta\) preserves \((f)+J\), and \(\delta\) preserves
\(J_A\). Multiplication maps also preserve \(J_A\). The endomorphisms
preserving a fixed ideal form a \(k\)-subalgebra, so they contain the entire
algebra generated by these maps. This covers every finite word and every
linear combination of words.
\end{proof}
The formal framework uses the additional hypothesis \(f\in J\), which is
proved explicitly for \(F\) by Euler's identity. It formalizes the polynomial
coefficient lift of every quotient derivation; no generation theorem for
Koszul derivations is assumed.

\section{A written homogeneous extension}\label{sec:general}
This section is a written proof, not a claim that the full class has been
formalized. Let \(k\) have characteristic zero, let \(n\ge3\), and let
\(f\in P=k[x_1,\ldots,x_n]\) be homogeneous of ordinary degree \(d\ge2\).
Write \(X=(x_i)^t\), \(H=\Hess(f)\), \(C=\adj(H)\), and \(h=\det H\).
Fix \(r\), put \(a=Ce_r\), and set
\begin{equation}\label{eq:generalQ}
 c=(n-2)(d-1),\qquad B=aX^t+Xa^t-x_rC,
 \qquad Q_r(p)=\sum_{i,j}B_{ij}\partial_i\partial_jp+c\sum_iC_{ir}\partial_ip.
\end{equation}

\begin{proposition}[Exact polynomial identities]\label{prop:identities}
For every \(p\in P\) and every \(j\),
\begin{align}
 Q_r(fp)&=f\left(Q_r(p)+2d\sum_i C_{ir}\partial_ip\right),\label{eq:descent}\\
 Q_r(f_j)&=(nd-n-d)\delta_{rj}h
       -(d-1)\sum_i(\partial_jC_{ir})f_i.\label{eq:escape}
\end{align}
\end{proposition}
\begin{proof}
Euler's identity and its derivative give
\[
 X^t\nabla f=df,\qquad HX=(d-1)\nabla f,
 \qquad C\nabla f=\frac{hX}{d-1}.
\]
Since \(H\) and \(C\) are symmetric,
\[
 B\nabla f=dfa,\qquad
 \tr(BH)=2a^tHX-x_r\tr(CH)=(2-n)x_rh.
\]
It follows that \(Q_r(f)=0\). The product rule and symmetry of \(B\) now yield
\[
 Q_r(fp)=fQ_r(p)+pQ_r(f)+2(B\nabla f)^t\nabla p,
\]
which proves \eqref{eq:descent} on every polynomial, not just on generators.

For \eqref{eq:escape}, the entries of \(H\) have degree \(d-2\), so
\(\sum_\ell x_\ell\partial_\ell H_{ij}=(d-2)H_{ij}\).
The two rank-one terms in \(B\), followed by the first-order part, contribute
\([2(d-2)+c]\delta_{rj}h\) to \(Q_r(f_j)\).
Jacobi's differentiation formula for the determinant gives the remaining term
\(-x_r\partial_jh\). On the other hand, differentiating
\((d-1)\sum_iC_{ir}f_i=x_rh\) gives
\[
 x_r\partial_jh=(d-2)\delta_{rj}h
       +(d-1)\sum_i(\partial_jC_{ir})f_i.
\]
Substitution proves the formula, since \(d-2+c=nd-n-d\).
\end{proof}

\begin{theorem}[Written homogeneous statement]\label{thm:general}
Assume additionally that \(P/J(f)\) is finite-dimensional. Each operator
\(Q_r\) in \eqref{eq:generalQ} descends to \(A=P/(f)\), has order at most
two, and lies outside \(\Delta_k(A)\).
\end{theorem}
\begin{proof}
Equation~\eqref{eq:descent} proves ideal preservation and quotient descent.
The displayed operator has differential order at most two; the triple
commutator identity passes to the quotient by surjectivity.

Euler's identity puts \(f\) in \(J(f)\). The gradient generators have positive
degree and define a zero-dimensional ideal, so they form a regular sequence
and \(P/J(f)\) is a nonzero Artinian local complete intersection supported
at the origin. It is already complete and has presentation
\(k[[x_1,\ldots,x_n]]/(f_1,\ldots,f_n)\).
The classical complete-intersection Jacobian--socle theorem, in the form of
Eisenbud--Ulrich \cite[Theorem 7.1, p.~210]{EU}, says that its K\"ahler
different is nonzero and equals its socle. Here the different is generated
by \(\det(\partial_jf_i)=h\). Thus \(h\notin J(f)\).

The integer \(nd-n-d=(n-1)(d-1)-1\) is positive under the stated hypotheses.
Equation~\eqref{eq:escape} shows that \(Q_r\) sends
\(\overline{f_r}\in J_A\) outside \(J_A\). Lemma~\ref{lem:stable}
therefore excludes \(Q_r\) from the full generated algebra.
\end{proof}
\begin{remark}
The isolatedness hypothesis is used for the classical socle conclusion, not
for the exact operator identities or derivation stability. Ordinary
homogeneity is essential to the argument given here; no weighted-homogeneous
extension is inferred by changing terminology.
\end{remark}

\section{A finite certificate over the rationals}\label{sec:finite}
For the fixed quartic \eqref{eq:F}, put \(S=\sum_i x_i^2\). Direct calculation
gives
\[
 G_i=2x_i(S+x_i^2),\qquad
 H_{ii}=2S+10x_i^2,\qquad H_{ij}=4x_ix_j\ (i\ne j).
\]
The specialized identities, checked in Lean for the actual Hessian adjugate,
are
\begin{align}
 \mathcal Q(Fp)&=F\left(\mathcal Q(p)+8\sum_i a_i\partial_ip\right),\label{eq:concretedescent}\\
 \mathcal Q(G_0)&=8h-3\sum_i(\partial_0a_i)G_i.\label{eq:concreteescape}
\end{align}
Define a \(\QQ\)-linear functional \(\lambda:P\to\QQ\) by its values on
monomials. It vanishes except in degree eight on the following sorted
exponent types:
\begin{center}
\begin{tabular}{lrrr}
\toprule
Exponent type & \(\lambda(x^\alpha)\) & coefficient of \(h\) & multiplicity\\
\midrule
\((8,0,0,0)\) & 3 & 96 & 4\\
\((6,2,0,0)\) & -2 & 720 & 12\\
\((4,4,0,0)\) & -2 & 1248 & 6\\
\((4,2,2,0)\) & 3 & 4128 & 12\\
\((2,2,2,2)\) & -12 & 16272 & 1\\
\bottomrule
\end{tabular}
\end{center}
\begin{lemma}\label{lem:finite}
For all \(i\) and all \(p\in P\), \(\lambda(G_ip)=0\), whereas
\(\lambda(h)=-77760\).
\end{lemma}
\begin{proof}
By linearity and degree, the annihilation calculation reduces to degree-five
monomials \(p\). Only those odd in \(x_i\) and even in the other variables
can contribute. Their types are \(x_i^5\), \(x_i^3x_j^2\),
\(x_ix_j^4\), and \(x_ix_j^2x_k^2\), with distinct indicated indices.
The corresponding evaluations are
\[
 4(3)+6(-2)=0,\quad
 4(-2)+2(-2)+4(3)=0,\quad
 4(-2)+2(-2)+4(3)=0,
\]
\[
 4(3)+2(3)+2(3)+2(-12)=0.
\]
All other parity patterns vanish termwise. Expanding the determinant gives
the coefficients in the table; consequently
\[
 \lambda(h)=4\cdot3\cdot96-12\cdot2\cdot720-6\cdot2\cdot1248
            +12\cdot3\cdot4128-12\cdot16272=-77760.
\]
The formal proof instead uses the coefficient formula for a monomial times
an arbitrary polynomial, proving \(\lambda(G_ip)=0\) directly for every
\(p\). The full 35-monomial determinant expansion and all functional values
are proved by kernel-checked algebraic normalization.
\end{proof}

\begin{proof}[Proof of Theorem~\ref{thm:fixed}]
Equation~\eqref{eq:concretedescent} gives an actual linear operator on the
quotient. Its triple commutators vanish because \(\mathcal Q\) is a sum of
polynomial coefficients times partial-derivative compositions of length at
most two. The quotient map transports this identity.
Euler's identity gives \(F\in J\), and Lemma~\ref{lem:finite} gives
\(h\notin J\). More explicitly, \eqref{eq:concreteescape} yields
\[
 \lambda(\mathcal Q(G_0))=8\lambda(h)=-622080\ne0,
\]
so \(\mathcal Q(G_0)\notin J\). As \((F)\subseteq J\), this remains an
escape from the image of \(J\) in the actual quotient. Apply
Lemma~\ref{lem:stable}. The concrete proof requires no socle theorem.
\end{proof}

\section{Geometric checks outside the formal theorem}
These observations are written proofs and are not asserted as additional
Lean declarations. They explain the example's geometric setting.
For \(F_n=\sum_i x_i^4+\sum_{i<j}x_i^2x_j^2\), \(n\ge3\), a common zero
of the gradient over an algebraic closure satisfies
\(x_i(S+x_i^2)=0\). If its nonzero support has size \(r\), then
\(x_i^2=-S\) on that support, whence \(S=-rS\). Characteristic zero gives
\(S=0\), and every coordinate vanishes. Thus the critical point is isolated.

There is no nontrivial linear Thom--Sebastiani splitting. Such a splitting
would produce a nontrivial constant idempotent matrix \(M\) with
\(H(X)M=M^tH(X)\). Comparing coefficients of \(x_k^2\) gives
\[
 (2+10\delta_{ik})M_{ij}=(2+10\delta_{jk})M_{ji}.
\]
For \(i\ne j\), choose first \(k\) distinct from both and then \(k=i\).
This forces \(M_{ij}=M_{ji}=0\). Comparing \(x_ix_j\) then equates all
diagonal entries. Thus \(M\) is scalar and has no nontrivial idempotent.
This excludes linear separated-variable forms; it is not a complete
analytic/contact-equivalence classification or a priority argument.

\section{Formal specification and reproducibility}
The formalization uses a four-variable multivariate polynomial algebra over
\(\QQ\), the actual ideal
quotient, \texttt{Derivation}, and \texttt{Algebra.adjoin}. The generated
algebra includes every multiplication map and every derivation of that
quotient; no hand-picked list of derivations substitutes for it.
The generic assembly theorem proves that four polynomial certificates imply
full quotient separation: \(f\in J(f)\), preservation of every multiple of
\(f\), the triple-commutator identity on the polynomial ring, and an actual
gradient-generator escape. All four are discharged for \eqref{eq:F}.

The checked source closure has fifteen modules. The final declarations report
only the standard axioms \texttt{propext}, \texttt{Classical.choice}, and
\texttt{Quot.sound}; there are no custom axioms, unfinished proof placeholders,
or native computation oracles. The pinned versions are:
\begin{lstlisting}
Lean 4.32.2: f3b06c705e6c85f5314019d5d3baab0fec5b580c
mathlib:     905b95818eb32af7874a58b427f50c1711a5e96c
\end{lstlisting}
The accompanying manifest records source hashes, the dependency order,
successful checks, and the independent replay. A clean rebuild of all fifteen
modules and the expanded final theorem also succeeded, using matching copied
source snapshots and a fresh output directory without the original object
files. Its log is included in the package. The supplied build
script uses one Lean thread and fixed time and memory guards.

Theorem~\ref{thm:fixed} is the full formal result. Theorem~\ref{thm:general}
remains a written theorem with the classical import stated explicitly.
No unrestricted Nakai result and no mathematical publication priority is
claimed.

\begin{thebibliography}{9}
\bibitem{DJM}
R.~N. Diethorn, J. Jeffries, C. Miller, N. Packauskas, J. Pollitz,
H. Rahmati, and S. Vassiliadou,
\emph{Resolutions of differential operators of low order for an isolated
hypersurface singularity}, arXiv:2209.13110v3, 16 November 2023.
\url{https://arxiv.org/html/2209.13110v3}.

\bibitem{YZZ}
S.~S.-T. Yau, Q. Zhu, and H. Zuo,
\emph{Nakai conjecture for isolated weighted homogeneous hypersurface
singularities}, arXiv:2604.24508v4, 8 October 2026.
\url{https://arxiv.org/html/2604.24508v4}.

\bibitem{BV}
P.~R. Brumatti and M.~O. Veloso,
\emph{A note on Nakai's conjecture for the ring
$K[X_1,\ldots,X_n]/(a_1X_1^m+\cdots+a_nX_n^m)$},
Colloquium Mathematicum \textbf{123} (2011), 277--283.
\url{https://doi.org/10.4064/cm123-2-10}.

\bibitem{EU}
D. Eisenbud and B. Ulrich,
\emph{Duality and socle generators for residual intersections},
Journal f\"ur die reine und angewandte Mathematik \textbf{756} (2019),
183--226, especially Theorem 7.1, p.~210.
\url{https://par.nsf.gov/servlets/purl/10282767}.

\bibitem{Traves}
W.~N. Traves,
\emph{Nakai's conjecture for varieties smoothed by normalization},
Proceedings of the American Mathematical Society \textbf{127} (1999),
2245--2248.
\url{https://www.usna.edu/Users/math/traves/_files/documents/nakaisconjecture.pdf}.
\end{thebibliography}
\end{document}
